
\documentclass[11pt]{article}
\usepackage[utf8]{inputenc}
%\usepackage[T1]{fontenc}
\usepackage{amssymb}
\usepackage{amsmath}
\usepackage{enumerate}
\usepackage{fullpage}
\usepackage{polski}  
\usepackage{indentfirst} 
\usepackage[pdftex]{graphicx}
\usepackage{multirow}
\usepackage{placeins}
\usepackage{wasysym}

\begin{document}

\begin{center}
\begin{tabular}{|l|l|l|}
\hline
	Wydział:EAIE Kierunek:EiT & & 
	\begin{tabular}{l}
		Rok I (2011/2012\\
		Grupa 6\\
		Zespół D\\
	\end{tabular}
 \\
\hline
	Data wykonania 5 marca 2012 & 
	\begin{tabular}{l}
	Laboratorium Metrologi: Cw 1 i 2\\
	Wprowadzenie do obsługi przyrządów pomiarowych \\
	oraz analiza błędów i niepewności pomiarowych\\
\end{tabular}
  & \\
\hline
	Zaliczenie & Podpis prowadzącego & Uwagi\\
\hline
\end{tabular}
\end{center}

\section{Pomiar bezpośredni napięcia stałego milimetrem cyfrowym}
\subsection{Schemat}
\begin{center}
\includegraphics[scale=2.5]{1.pdf}
\end{center}
\subsection{Wyniki pomiarów}
\begin{center}
\begin{tabular}{|l|l|l|l|l|}
	\hline Lp. & U & $\Delta_{gr}U$ & $u_B(U)$ & $U(U)$ \\
	\hline 1 & 5,975 & 0,004 & 0,0023 & 0,004\\
	\hline 2 & 11,999 & 0,006 & 0,004 & 0,007\\
	\hline
\end{tabular}
\end{center}
\subsection{Wzory i przykładowe obliczenia}
Błąd graniczy bezwzględny pomiaru napięcia obliczamy z wzoru $$ \Delta_{gr} = \frac{a \cdot U + b \cdot Z_U}{100} $$

Przy pomiarze napięcia $U_1$ przy zakresie $40V$ współczynniki wynoszą $ a = 0,025 \quad b = 0,006 $ więc błąd graniczny wynosi 

$$ \Delta_{gr}U = \frac{0,025 \cdot 5,975 + 0,006 \cdot 40}{100} = 0,004$$

Niepewność typu B wynosi 
$$ u_B(U) = \frac{\Delta_{gr}U}{\sqrt{3}} = 0,0023$$

$$ U(U) = k u_B(U) \quad k = \sqrt{3}p \quad p = 0,95 $$
$$ U(U) = 0,004 $$

\subsection{Wnioski}
Napięcie dla zakresu $12V$ jest bardziej stabilne, niż dla $6Vo$. Napięcie na zaciskach zasilacza podczas gdy był ustawiony na $6V$ jest mniejsze, niż powinno być, a różnica jest większa, niż błąd graniczny.

\subsection{Wykaz urządzeń}
Rigol DM3051, Zasilacz laboratoryjny

\newpage
\section{Pomiar bezpośredni napięcia stałego multimetrem}

\subsection{Schemat}
\begin{center}
\includegraphics[]{2.pdf}
\end{center}
\subsection{Wyniki pomiarów}
\begin{center}
\begin{tabular}{|l|l|l|l|l|l|l|l|l|l|}
\hline
	LP & Klasa & $\alpha_m$ & $Z_{U}[V]$ & $\alpha[dz]$ & $c_U[V/dz]$ & $U[V]$ & $\Delta_{gr}[V]$ & $u_b(U)$ & $U(U)$\\
\hline
	1 & 1 & 60 & 6 & 59 & 0,1 & 5,9 & 0,06 & $\frac{\sqrt{3}}{50}$ & 0,06\\
\hline
	2 & 1 & 60 & 15 & 36 & 0,25 & 9 & 0,15 & $\frac{\sqrt{3}}{20}$ & 0,15\\
\hline
	3 & 1 & 60 & 30 & 24 & 0,5 & 12 & 0,3 & $\frac{\sqrt{3}}{10}$ & 0,3 \\
\hline
\end{tabular}
\end{center}
\subsection{Wzory i przykładowe obliczenia}
Obliczania przeprowadzone dla drugiego pomiaru
$$ \Delta_{gr}U = \frac{K \cdot Z_U}{100} = \frac{1 \cdot 15}{100} = 0,15 V$$

$$ u_b(U) = \frac{\Delta_{gr}U}{\sqrt{3}} = \frac{0.15}{\sqrt{3}} = 0,09 $$

$$ U(U) = k u_b(U) = 0,95 \sqrt{3} u_b(U) = 1,5 $$
\subsection{Wnioski}
Dla multimetru analogowego dokładność jest mniejsza. Nie dość, że błąd graniczny jest rząd wielkości większy, to pomiar jest obarczony dodatkowymi błędami (np. błąd paralaksy).
\subsection{Lista przyrządów}
UM3a, Zasilacz laboratoryjny

\newpage
\section{Pomiar bezpośredni rezystancji metodą dwuprzepływową} 
\subsection{Schemat}
\begin{center}
\includegraphics[scale=1.4]{3.pdf}
\end{center}
\subsection{Wyniki pomiarów}
\begin{center}
\begin{tabular}{|l|l|l|l|l|}
\hline
	 Rezystor & $R[k\Omega]$ & $Z_R[k\Omega]$ & $\Delta_{gr}R [\Omega]$ & $U(R)$\\
\hline
	R2 & 5,894 & $40$ & 25 & 24\\
\hline
	R4 & 36,526 & $40$ & 33 & 31\\
\hline
\end{tabular}
\end{center}
\subsection{Wzory i przykładowe obliczenia}
Obliczenia wykonywane dla $R2$
$$\Delta_{gr}R = \frac{a \cdot R + b \cdot Z_R}{100} = 
	\frac{5,894 *10^3 * 0,025 + 4*10^5*0,006}{100} = 25,674$$
Niepewność typu $B$ 
$$ U_B(R) = \frac{\Delta_{gr}R}{\sqrt{3}} = \frac{25,674}{\sqrt{3}} = 14,82 $$

$$ U(R) = \sqrt{3} \cdot p \cdot u_B(R) = 24,39 \quad p = 0,95$$

\subsection{Wnioski}
Pomiar znacznej rezystancji jest dokładny i błąd graniczny stanowi mniej niż 5$\permil$ całkowitego wyniku. 
\subsection{Wykaz przyrządów}
Płytka z rezystorami, Rigol DM3051

\section{Pomiar bezpośredni rezystancji metodą czteroprzewodową multimetrem cyforwym}

\subsection{Schemat}
\begin{center}
\includegraphics{4.pdf}
\end{center}
\subsection{Wyniki pomiarów}
\begin{center}
\begin{tabular}{|l|l|l|l|l|l|l|}
\hline
	Lp. & $R^{4W}$ & $\Delta_{gr}R^{4W}$ & $U(R^{4W})$ & $R^{2W}$ & $\Delta_{gr}R^{2W}$ & $U(R^{2W})$\\
\hline
	 1 & 1,048 & 0,004 & 0,0024 & 1,088 & 0,004 & 0,0024 \\
\hline
\end{tabular}
\end{center}
\subsection{Przykładowe obliczenia i wzory}
Błędy graniczne obliczamy z wzorów 

$$\Delta_{gr}R^{4W} = \frac{ a \cdot R + b \cdot Z_r}{100} [\Omega]$$ 
$$ \Delta_{gr}R^{2R} = \frac{ a \cdot R + b \cdot Z_r}{100} + 0,2 [\Omega]$$

Dla zakresu $200 \Omega$ i użytego multimetru $a = 0,010$ i $b = 0,0020$

$$ \Delta_{gr}R^{4W} = 0,4 $$
$$ \Delta_{gr}R^{2W} = 0,6 $$

\subsection{Wykaz przyrządów}
Płytka z rezystorem, Rigol DM3051

\newpage
\section{Pomiar bezpośredni napięcia skutecznego prądu zmiennego przy pomocy multimetru cyfrowego}
\subsection{Schemat}
\begin{center}
\includegraphics{5.pdf}
\end{center}
\subsection{Wyniki pomiarów}
\begin{center}
\begin{tabular}{|l|l|l|l|}
\hline
	Lp. & $U_{rms}[V]$ & $\Delta_{gr}U_{rms}[V]$ & $ U(U_{rms}) $\\
\hline
	1 & 9,719 & 0,006 & 0,004\\
\hline
	2 & 3,557 & 0,003 & 0,002\\
\hline
	3 & 8,484 & 0,005 & 0,003\\
\hline
\end{tabular}
\end{center}
\subsection{Przykładowe obliczenia i wzory}
$$ \Delta_{gr}U = \frac{a \cdot U + b \cdot Z_u}{100} [V] $$
Obliczenia przeprowadzamy na zakresie $20V$, wtedy $a = 0,04$ i $b = 0,01$. Obliczenia konkretne dla trzeciego pomiaru
$$ \Delta_{gr}U = \frac{8,484 \cdot 0,04 + 0,01 \cdot 20}{100} \approx 0,005 [V] $$

\subsection{Wnioski}
Pomiar jest dość dokładny, jednak przyrząd znacznie gorzej radzi sobie z napięciem zmiennym niż stałym.
\subsection{Wykaz urządzeń}
Zasilacz laboratoryjny, Rigol DM3051
\newpage
\section{Pomiar bezpośredni częstotliwości multimetrem cyfrowym}
\subsection{Schemat}

\begin{center}
\includegraphics{6.pdf}
\end{center}

\subsection{Wyniki pomiarów}

\begin{center}
\begin{tabular}{|l|l|l|l|}
\hline
	Lp. & $f[kHz]$ & $\Delta_{gr}f[Hz]$ & $U(f)[Hz]$\\
\hline
	1 & 1,231  & 0,4& 0,4\\
\hline
	2 & 2,127 & 0,4 & 0,4\\
\hline
\end{tabular}
\end{center}

\subsection{Wzory i przykładowe obliczenia}
$$ \Delta_{gr}f = \frac{a \cdot f}{100} [Hz] $$
Dla zakresu $3 - 5 kHz$ współczynnik $a = 0,02$. 
Obliczenia przeprowadzone dla drugiego pomiaru

$$ \Delta_{gr}f = \frac{0,02 \cdot 2,127}{100} = 0,0004  kHz$$   

\subsection{Wnioski}
Błędy graniczne stanowią niecały procent wyniku. Oznacza to, że pomiary częstotliwości są dość dokładne.
\subsection{Wykaz urządzeń}
Zasilacz laboratoryjny, Rigol DM3051

\newpage
\section{Pomiar częstotliwości napięcia zmiennego metodami z wykorzystaniem funkcji statystycznych}
\subsection{Schemat}
\begin{center}
\includegraphics{6.pdf}
\end{center}
\subsection{Wyniki pomiarów}
\begin{center}
\begin{tabular}{|l|l|l|l|l|}
\hline
	Lp. & n & $f_{min}$ & $f_{ave}$ & $f_{max}$ \\
\hline
	1 & 31 & 0,1373 & 0,1373 & 0,1374\\
\hline
	2 & 33 & 2,134 & 2,135 & 2,136\\
\hline
\end{tabular}
\end{center}
\subsection{Wnioski}
Obliczenie błędów i niepewności w tego typu pomiarach wykracza poza ramy ćwiczenia.
\subsection{Wykaz urządzeń}
Zasilacz laboratoryjny, Rigol DM3051

\newpage
\section{Bezpośredni pomiar pojemności z użyciem metod statystycznych}
\subsection{Schemat}
\begin{center}
\includegraphics[scale=1.3]{8.pdf}
\end{center}
\subsection{Wyniki pomiarów}
\begin{center}
\begin{tabular}{|l|l|l|l|l|}
\hline
	Lp. & n & $C_{min}[nF]$ & $C_{ave}[nF]$ & $C_{max}[nF]$\\
\hline
	1 & 36 & 2,1074 & 2,1131 & 2,1296\\
\hline
	2 & 31 & 29,794 & 29,808 & 29,843\\
\hline
\end{tabular}
\end{center}
\subsection{Wnioski}
Wyznaczenie błędu oraz niepewności w tego typu pomiarach wykracza poza ramy ćwiczenia.
\subsection{Wykaz urządzeń}
Zasilacz laboratoryjny, Rigol DM3051
\end{document}
